\chapter{2014年新课标II卷（文）}

\section{单选题}

\begin{problem}
已知集合 \(A=\{-2,0,2\}\)，\(B=\{x\mid x^2-x-2=0\}\)，则 \(A\cap B=\)（\quad）

\choices
    {\(\varnothing\)}
    {\(\{2\}\)}
    {\(\{0\}\)}
    {\(\{-2\}\)}
\end{problem}

\begin{answer}
B
\end{answer}

\begin{solution}
由 \(x^2-x-2=(x-2)(x+1)\)，得 \(B=\{-1,2\}\). 因此 \(A\cap B=\{2\}\)，故选 B.
\end{solution}

\begin{problem}
\(\frac{1+3\i}{1-\i}=\)（\quad）

\choices
    {\(1+2\i\)}
    {\(-1+2\i\)}
    {\(1-2\i\)}
    {\(-1-2\i\)}
\end{problem}

\begin{answer}
B
\end{answer}

\begin{solution}
分子、分母同乘分母的共轭复数 \(1+\i\)，得
\(\displaystyle\frac{1+3\i}{1-\i}=\frac{(1+3\i)(1+\i)}{(1-\i)(1+\i)}=\frac{1+\i+3\i+3\i^2}{1-\i^2}=\frac{-2+4\i}{2}=-1+2\i\). 故选 B.
\end{solution}

\begin{problem}
函数 \(f(x)\) 在 \(x=x_0\) 处导数存在，若 \(p\)：\(f'(x_0)=0\)；\(q\)：\(x=x_0\) 是 \(f(x)\) 的极值点，则（\quad）

\choices
    {\(p\) 是 \(q\) 的充分必要条件}
    {\(p\) 是 \(q\) 的充分条件，但不是 \(q\) 的必要条件}
    {\(p\) 是 \(q\) 的必要条件，但不是 \(q\) 的充分条件}
    {\(p\) 既不是 \(q\) 的充分条件，也不是 \(q\) 的必要条件}
\end{problem}

\begin{answer}
C
\end{answer}

\begin{solution}
若 \(x=x_0\) 是可导函数的极值点，由费马定理知 \(f'(x_0)=0\)，所以 \(q\) 是 \(p\) 的充分条件，即 \(p\) 是 \(q\) 的必要条件. 反之不成立，例如 \(f(x)=(x-x_0)^3\) 满足 \(f'(x_0)=0\)，但 \(x=x_0\) 不是极值点. 因此 \(p\) 是 \(q\) 的必要条件，但不是充分条件，故选 C.
\end{solution}

\begin{problem}
设向量 \(\bs{a}\)，\(\bs{b}\) 满足 \(|\bs{a}+\bs{b}|=\sqrt{10}\)，\(|\bs{a}-\bs{b}|=\sqrt{6}\)，则 \(\bs{a}\cdot\bs{b}=\)（\quad）

\choices
    {\(1\)}
    {\(2\)}
    {\(3\)}
    {\(5\)}
\end{problem}

\begin{answer}
A
\end{answer}

\begin{solution}
将两个已知等式平方，得
\(|\bs{a}+\bs{b}|^2=|\bs{a}|^2+2\bs{a}\cdot\bs{b}+|\bs{b}|^2=10\).

\(|\bs{a}-\bs{b}|^2=|\bs{a}|^2-2\bs{a}\cdot\bs{b}+|\bs{b}|^2=6\).
两式相减得 \(4\bs{a}\cdot\bs{b}=4\)，所以 \(\bs{a}\cdot\bs{b}=1\)，故选 A.
\end{solution}

\begin{problem}
等差数列 \(\{a_n\}\) 的公差是 \(2\)，若 \(a_2\)，\(a_4\)，\(a_8\) 成等比数列，则 \(\{a_n\}\) 的前 \(n\) 项和 \(S_n=\)（\quad）

\choices
    {\(n(n+1)\)}
    {\(n(n-1)\)}
    {\(\frac{n(n+1)}{2}\)}
    {\(\frac{n(n-1)}{2}\)}
\end{problem}

\begin{answer}
A
\end{answer}

\begin{solution}
设 \(a_4=u\). 由于公差为 \(2\)，有 \(a_2=u-4\)，\(a_8=u+8\). 三项 \(a_2\)，\(a_4\)，\(a_8\) 成等比数列，所以中间项的平方等于两边两项的积，即
\(u^2=(u-4)(u+8)=u^2+4u-32\)，从而 \(u=8\). 因而 \(a_1=a_4-3\times2=2\)，且 \(a_n=a_1+(n-1)\times2=2n\). 所以
\(S_n=\frac{n(a_1+a_n)}{2}=\frac{n(2+2n)}{2}=n(n+1)\)，故选 A.
\end{solution}

\begin{problem}
如图，网格纸上正方形小格的边长为 \(1\)（表示 \(1\mathrm{~cm}\)），图中粗线画出的是某零件的三视图，该零件由一个底面半径为 \(3\mathrm{~cm}\)，高为 \(6\mathrm{~cm}\) 的圆柱体毛坯切削得到，则切削掉的部分的体积与原来毛坯体积的比值为（\quad）

\bitmapfigure[width=0.42\linewidth]{img/2014/new_curriculum_paper_2_liberal/q6_fig1.png}

\choices
    {\(\frac{17}{27}\)}
    {\(\frac{5}{9}\)}
    {\(\frac{10}{27}\)}
    {\(\frac{1}{3}\)}
\end{problem}

\begin{answer}
C
\end{answer}

\begin{solution}
由俯视图中两个同心圆的半径和正、侧视图中的高度可知，零件上部是底面半径为 \(2\mathrm{~cm}\)、高为 \(4\mathrm{~cm}\) 的圆柱，下部是底面半径为 \(3\mathrm{~cm}\)、高为 \(2\mathrm{~cm}\) 的圆柱. 因此保留零件的体积为
\(V_1=\pi\cdot2^2\cdot4+\pi\cdot3^2\cdot2=34\pi\mathrm{~cm}^3\).
原毛坯体积为 \(V_0=\pi\cdot3^2\cdot6=54\pi\mathrm{~cm}^3\)，切削掉的体积为 \(V_0-V_1=20\pi\mathrm{~cm}^3\). 所求比值为
\(\displaystyle\frac{V_0-V_1}{V_0}=\frac{20\pi}{54\pi}=\frac{10}{27}\)，故选 C.
\end{solution}

\begin{problem}
正三棱柱 \(ABC-A_1B_1C_1\) 的底面边长为 \(2\)，侧棱长为 \(\sqrt{3}\)，\(D\) 为 \(BC\) 中点，则三棱锥 \(A-B_1DC_1\) 的体积为（\quad）

\choices
    {\(3\)}
    {\(\frac{3}{2}\)}
    {\(1\)}
    {\(\frac{\sqrt{3}}{2}\)}
\end{problem}

\begin{answer}
C
\end{answer}

\begin{solution}
建立空间直角坐标系，取 \(D=(0,0,0)\)，\(B=(-1,0,0)\)，\(C=(1,0,0)\)，\(A=(0,\sqrt3,0)\)，并令侧棱方向为 \(z\) 轴正方向. 于是 \(B_1=(-1,0,\sqrt3)\)，\(C_1=(1,0,\sqrt3)\).
四面体 \(A-B_1DC_1\) 的体积为
\(\displaystyle\frac{1}{6}\left|\overrightarrow{DA}\cdot\left(\overrightarrow{DB_1}\times\overrightarrow{DC_1}\right)\right|\).
其中 \(\overrightarrow{DB_1}=(-1,0,\sqrt3)\)，\(\overrightarrow{DC_1}=(1,0,\sqrt3)\)，所以 \(\overrightarrow{DB_1}\times\overrightarrow{DC_1}=(0,2\sqrt3,0)\)，从而体积为 \(\frac{1}{6}(2\sqrt3\cdot\sqrt3)=1\)，故选 C.
\end{solution}

\begin{problem}
执行如图所示的程序框图，如果输入的 \(x\)，\(t\) 均为 \(2\)，则输出的 \(S=\)（\quad）

\bitmapfigure[width=0.46\linewidth]{img/2014/new_curriculum_paper_2_liberal/q8_fig1.png}
\FloatBarrier

\choices
    {\(4\)}
    {\(5\)}
    {\(6\)}
    {\(7\)}
\end{problem}

\begin{answer}
D
\end{answer}

\begin{solution}
开始时 \(M=1\)，\(S=3\)，\(k=1\). 第一次循环满足 \(k\leqslant t\)，于是 \(M=\frac{1}{1}\cdot2=2\)，\(S=2+3=5\)，\(k=2\). 第二次循环中 \(M=\frac{2}{2}\cdot2=2\)，\(S=2+5=7\)，\(k=3\). 此时 \(k>t=2\)，循环结束，输出 \(S=7\)，故选 D.
\end{solution}

\begin{problem}
设 \(x\)，\(y\) 满足约束条件 \(\begin{cases}
x+y-1\geqslant 0,\\
x-y-1\leqslant 0,\\
x-3y+3\geqslant 0,
\end{cases}\) 则 \(z=x+2y\) 的最大值为（\quad）

\choices
    {\(8\)}
    {\(7\)}
    {\(2\)}
    {\(1\)}
\end{problem}

\begin{answer}
B
\end{answer}

\begin{solution}
约束条件等价于 \(y\geqslant1-x\)，\(y\geqslant x-1\)，\(y\leqslant\frac{x}{3}+1\). 三条边界直线的交点分别由
\(1-x=x-1\)、\(x-1=\frac{x}{3}+1\)、\(1-x=\frac{x}{3}+1\) 解得 \((1,0)\)、\((3,2)\)、\((0,1)\)，这三个点都满足另外一个不等式，因此它们正是可行域的顶点.
在这三个顶点处，\(z=x+2y\) 的值分别为 \(1\)、\(7\)、\(2\). 线性目标函数在有界多边形可行域上的最大值取在顶点处，所以最大值为 \(7\)，故选 B.
\end{solution}

\begin{problem}
设 \(F\) 为抛物线 \(C:y^2=3x\) 的焦点，过 \(F\) 且倾斜角为 \(30^{\circ}\) 的直线交 \(C\) 于 \(A\)，\(B\) 两点，则 \(|AB|=\)（\quad）

\choices
    {\(\frac{\sqrt{30}}{3}\)}
    {\(6\)}
    {\(12\)}
    {\(7\sqrt{3}\)}
\end{problem}

\begin{answer}
C
\end{answer}

\begin{solution}
将抛物线写成 \(y^2=4px\)，得 \(p=\frac34\)，焦点为 \(F\left(\frac34,0\right)\). 用参数方程 \(x=pt^2\)，\(y=2pt\) 表示抛物线上的点. 设 \(A\)、\(B\) 的参数分别为 \(t_1\)，\(t_2\)，则
\(\displaystyle k_{AB}=\frac{2p(t_1-t_2)}{p(t_1^2-t_2^2)}=\frac{2}{t_1+t_2}\).
由两参数点确定的弦方程 \(y(t_1+t_2)=2x+2pt_1t_2\)，将焦点 \(F\left(p,0\right)\) 代入，得 \(t_1t_2=-1\).
直线倾斜角为 \(30^{\circ}\)，所以
\(\displaystyle\frac{2}{t_1+t_2}=\tan30^{\circ}=\frac{1}{\sqrt3}\)，从而 \(t_1+t_2=2\sqrt3\). 于是
\(\lvert t_1-t_2\rvert=\sqrt{(t_1+t_2)^2-4t_1t_2}=\sqrt{12+4}=4\).
两点间距离为
\(\displaystyle|AB|=p|t_1-t_2|\sqrt{(t_1+t_2)^2+4}=\frac34\cdot4\cdot4=12\)，故选 C.
\end{solution}

\begin{problem}
若函数 \(f(x)=kx-\ln x\) 在区间 \((1,+\infty)\) 单调递增，则 \(k\) 的取值范围是（\quad）

\choices
    {\((-\infty,-2]\)}
    {\((-\infty,-1]\)}
    {\([2,+\infty)\)}
    {\([1,+\infty)\)}
\end{problem}

\begin{answer}
D
\end{answer}

\begin{solution}
函数在 \((1,+\infty)\) 上的导数为 \(f'(x)=k-\frac1x\). 若 \(k\geqslant1\)，则对任意 \(x>1\) 有 \(\frac1x<1\leqslant k\)，从而 \(f'(x)>0\)，所以函数单调递增.
若 \(k\leqslant0\)，任取 \(x>1\) 都有 \(f'(x)=k-\frac1x<0\)，不可能在整个区间上单调递增；若 \(0<k<1\)，取 \(x\) 满足 \(1<x<\frac1k\)，则 \(\frac1x>k\)，仍有 \(f'(x)<0\). 因此 \(k\geqslant1\)，故选 D.
\end{solution}

\begin{problem}
设点 \(M(x_0,1)\)，若在圆 \(O:x^2+y^2=1\) 上存在点 \(N\)，使得 \(\angle OMN=45^{\circ}\)，则 \(x_0\) 的取值范围是（\quad）

\choices
    {\([-1,1]\)}
    {\(\left[-\frac{1}{2},\frac{1}{2}\right]\)}
    {\(\left[-\sqrt{2},\sqrt{2}\right]\)}
    {\(\left[-\frac{\sqrt{2}}{2},\frac{\sqrt{2}}{2}\right]\)}
\end{problem}

\begin{answer}
A
\end{answer}

\begin{solution}
设 \(M=(x_0,1)\)，记 \(r=OM=\sqrt{x_0^2+1}\)，并令 \(s=MN>0\). 在三角形 \(OMN\) 中，\(ON=1\)，且 \(\angle OMN=45^{\circ}\)，由余弦定理得
\(1=r^2+s^2-2rs\cos45^{\circ}=r^2+s^2-\sqrt2rs\).
把它看作关于 \(s\) 的二次方程 \(s^2-\sqrt2rs+r^2-1=0\)，其判别式为
\(\Delta=2r^2-4(r^2-1)=4-2r^2\). 因而存在实数 \(s\) 的必要条件是 \(r^2\leqslant2\)，即 \(x_0^2\leqslant1\).
反之，若 \(x_0^2\leqslant1\)，则 \(1\leqslant r\leqslant\sqrt2\)，取方程的正根
\(\displaystyle s=\frac{\sqrt2r+\sqrt{4-2r^2}}{2}>0\).
从 \(M\) 出发作与 \(MO\) 成 \(45^{\circ}\) 的射线，在该射线上取距 \(M\) 为 \(s\) 的点 \(N\). 上式保证 \(ON=1\)，所以 \(N\) 在圆 \(O\) 上且 \(\angle OMN=45^{\circ}\). 因此 \(x_0\in[-1,1]\)，故选 A.
\end{solution}

\section{填空题}

\begin{problem}
甲，乙两名运动员各自等可能地从红，白，蓝 \(3\) 种颜色的运动服中选择 \(1\) 种，则他们选择相同颜色运动服的概率为\fillinblank{}.
\end{problem}

\begin{answer}
\(\frac{1}{3}\)
\end{answer}

\begin{solution}
甲运动员有 \(3\) 种等可能选择，乙运动员也有 \(3\) 种等可能选择，所以所有等可能结果共有 \(3\times3=9\) 个. 两人选同色的结果为红红、白白、蓝蓝，共 \(3\) 个，故所求概率为 \(\frac{3}{9}=\frac13\).
\end{solution}

\begin{problem}
函数 \(f(x)=\sin\left(x+\varphi\right)-2\sin\varphi\cos x\) 的最大值为\fillinblank{}.
\end{problem}

\begin{answer}
1
\end{answer}

\begin{solution}
利用和角公式，得 \(f(x)=\sin x\cos\varphi+\cos x\sin\varphi-2\sin\varphi\cos x=\sin x\cos\varphi-\cos x\sin\varphi=\sin(x-\varphi)\). 正弦函数的最大值为 \(1\)，所以 \(f(x)\) 的最大值为 \(1\).
\end{solution}

\begin{problem}
偶函数 \(y=f(x)\) 的图象关于直线 \(x=2\) 对称，\(f(3)=3\)，则 \(f(-1)=\)\fillinblank{}.
\end{problem}

\begin{answer}
3
\end{answer}

\begin{solution}
函数为偶函数，所以 \(f(-1)=f(1)\). 图象关于直线 \(x=2\) 对称，反射 \(x=3\) 得到 \(x=1\)，故 \(f(1)=f(3)=3\). 因此 \(f(-1)=3\).
\end{solution}

\begin{problem}
数列 \(\{a_n\}\) 满足 \(a_{n+1}=\frac{1}{1-a_n}\)，\(a_8=2\)，则 \(a_1=\)\fillinblank{}.
\end{problem}

\begin{answer}
\(\frac{1}{2}\)
\end{answer}

\begin{solution}
由递推式反解得 \(a_n=1-\frac{1}{a_{n+1}}\). 从 \(a_8=2\) 逐项向前计算：\(a_7=\frac12\)，\(a_6=-1\)，\(a_5=2\)，\(a_4=\frac12\)，\(a_3=-1\)，\(a_2=2\)，\(a_1=\frac12\). 因此 \(a_1=\frac12\).
\end{solution}

\section{解答题}

\begin{problem}
四边形 \(ABCD\) 的内角 \(A\) 与 \(C\) 互补，\(AB=1\)，\(BC=3\)，\(CD=DA=2\).
\begin{enumerate}
\item 求 \(\angle C\) 和 \(BD\)；
\item 求四边形 \(ABCD\) 的面积.
\end{enumerate}
\end{problem}

\begin{solution}
\begin{enumerate}
\item 记 \(\angle C=\gamma\). 因为 \(\angle A+\angle C=180^{\circ}\)，所以四边形 \(ABCD\) 为圆内接四边形，且 \(\angle A=180^{\circ}-\gamma\). 在三角形 \(BCD\) 和 \(BAD\) 中分别使用余弦定理，得
\(BD^2=BC^2+CD^2-2BC\cdot CD\cos\gamma=13-12\cos\gamma\)，
\(BD^2=AB^2+AD^2-2AB\cdot AD\cos(180^{\circ}-\gamma)=5+4\cos\gamma\).
两式相等，得 \(13-12\cos\gamma=5+4\cos\gamma\)，所以 \(\cos\gamma=\frac12\). 因为 \(0^{\circ}<\gamma<180^{\circ}\)，故 \(\gamma=60^{\circ}\)，且 \(BD^2=13-12\cdot\frac12=7\)，从而 \(BD=\sqrt7\).
\item 四边形面积等于两个三角形面积之和，因此
\(S_{ABCD}=\frac12 AB\cdot AD\sin\angle A+\frac12 BC\cdot CD\sin\angle C\).
又 \(\sin\angle A=\sin\angle C=\sin60^{\circ}=\frac{\sqrt3}{2}\)，所以 \(S_{ABCD}=\sin60^{\circ}+3\sin60^{\circ}=2\sqrt3\).
\end{enumerate}
\end{solution}

\begin{problem}
如图，四棱锥 \(P-ABCD\) 中，底面 \(ABCD\) 为矩形，\(PA\perp\) 平面 \(ABCD\)，\(E\) 为 \(PD\) 的中点.
\begin{enumerate}
\item 证明：\(PB\parallel\) 平面 \(AEC\)；
\item 设 \(AP=1\)，\(AD=\sqrt{3}\)，三棱锥 \(P-ABD\) 的体积 \(V=\frac{\sqrt{3}}{4}\)，求 \(A\) 到平面 \(PBC\) 的距离.
\end{enumerate}

\bitmapfigure[width=0.38\linewidth]{img/2014/new_curriculum_paper_2_liberal/q18_fig1.png}
\FloatBarrier
\end{problem}

\begin{solution}
\begin{enumerate}
\item 设 \(AB=b\)，\(AD=d\)，\(AP=h\)，建立直角坐标系，使
\(A=(0,0,0)\)，\(B=(b,0,0)\)，\(D=(0,d,0)\)，\(C=(b,d,0)\)，\(P=(0,0,h)\). 因为 \(E\) 是 \(PD\) 的中点，所以 \(E=(0,\frac d2,\frac h2)\).
于是 \(\overrightarrow{PB}=(b,0,-h)\)，而平面 \(AEC\) 的一个法向量为
\(\bs{n}=\overrightarrow{AC}\times\overrightarrow{AE}=(\frac{dh}{2},-\frac{bh}{2},\frac{bd}{2})\).
计算得 \(\overrightarrow{PB}\cdot\bs{n}=\frac{bdh}{2}-\frac{bdh}{2}=0\). 又 \(\overrightarrow{AB}\cdot\bs{n}=\frac{bdh}{2}\ne0\)，所以 \(B\notin\) 平面 \(AEC\)，从而 \(PB\parallel\) 平面 \(AEC\).
\item 三棱锥 \(P-ABD\) 的底面三角形 \(ABD\) 的面积为 \(\frac12bd\)，故
\(V=\frac13\cdot\frac12bd\cdot h=\frac{bdh}{6}\). 代入 \(h=1\)，\(d=\sqrt3\)，\(V=\frac{\sqrt3}{4}\)，得 \(b=\frac32\).
此时 \(P=(0,0,1)\)，\(B=(\frac32,0,0)\)，\(C=(\frac32,\sqrt3,0)\). 平面 \(PBC\) 的一个法向量可取 \((2,0,3)\)，其方程为 \(2x+3z-3=0\). 因而点 \(A=(0,0,0)\) 到该平面的距离为
\(\displaystyle\frac{|2\cdot0+3\cdot0-3|}{\sqrt{2^2+3^2}}=\frac{3}{\sqrt{13}}=\frac{3\sqrt{13}}{13}\).
\end{enumerate}
\end{solution}

\begin{problem}
某市为了考核甲，乙两部门的工作情况，随机访问了 \(50\) 位市民，根据这 \(50\) 位市民对这两部门的评分（评分越高表明市民的评价越高），绘制茎叶图如下：

\begin{enumerate}
\item 分别估计该市的市民对甲，乙两部门评分的中位数；
\item 分别估计该市的市民对甲，乙两部门的评分高于 \(90\) 的概率；
\item 根据茎叶图分析该市的市民对甲，乙两部门的评价.
\end{enumerate}

\begin{center}
\renewcommand{\arraystretch}{1.3}
\begin{tabular}{r|c|l}
\multicolumn{1}{c|}{甲部门} & \multicolumn{1}{c|}{} & \multicolumn{1}{c}{乙部门} \\
\hline
 & \(3\) & \(5\enspace 9\) \\
\(4\) & \(4\) & \(0\enspace 4\enspace 4\enspace 8\) \\
\(9\enspace 7\) & \(5\) & \(1\enspace 2\enspace 2\enspace 4\enspace 5\enspace 6\enspace 6\enspace 7\enspace 7\enspace 7\enspace 8\enspace 9\) \\
\(9\enspace 7\enspace 6\enspace 6\enspace 5\enspace 3\enspace 3\enspace 2\enspace 1\enspace 1\enspace 0\) & \(6\) & \(0\enspace 1\enspace 1\enspace 2\enspace 3\enspace 4\enspace 6\enspace 8\enspace 8\) \\
\(9\enspace 8\enspace 8\enspace 7\enspace 7\enspace 7\enspace 6\enspace 6\enspace 5\enspace 5\enspace 5\enspace 5\enspace 5\enspace 4\enspace 4\enspace 4\enspace 3\enspace 3\enspace 3\enspace 2\enspace 1\enspace 0\enspace 0\) & \(7\) & \(0\enspace 0\enspace 1\enspace 1\enspace 3\enspace 4\enspace 4\enspace 9\) \\
\(6\enspace 6\enspace 5\enspace 5\enspace 2\enspace 0\enspace 0\) & \(8\) & \(1\enspace 2\enspace 3\enspace 3\enspace 4\enspace 5\) \\
\(6\enspace 3\enspace 2\enspace 2\enspace 2\enspace 0\) & \(9\) & \(0\enspace 1\enspace 1\enspace 4\enspace 5\enspace 6\) \\
 & \(10\) & \(0\enspace 0\enspace 0\)
\end{tabular}
\end{center}
\end{problem}

\begin{solution}
\begin{enumerate}
\item 甲部门的 \(50\) 个评分中，前 \(14\) 个不超过 \(69\)，第 \(15\) 至第 \(37\) 个在 \(70\) 至 \(79\) 之间. 第 \(25\)、\(26\) 个评分在该行叶子序列中分别都是 \(5\)，所以甲部门评分的中位数为 \(\frac{75+75}{2}=75\). 乙部门前 \(18\) 个评分不超过 \(59\)，第 \(19\) 至第 \(27\) 个在 \(60\) 至 \(68\) 之间，第 \(25\)、\(26\) 个评分分别为 \(66\)、\(68\)，所以乙部门评分的中位数为 \(\frac{66+68}{2}=67\).
\item 甲部门评分中高于 \(90\) 的有 \(92\)，\(92\)，\(92\)，\(93\)，\(96\)，共 \(5\) 个，故概率估计为 \(\frac5{50}=0.1\). 乙部门评分中高于 \(90\) 的有 \(91\)，\(91\)，\(94\)，\(95\)，\(96\)，\(100\)，\(100\)，\(100\)，共 \(8\) 个，故概率估计为 \(\frac8{50}=0.16\).
\item 甲部门评分的中位数高于乙部门，且评分主要集中在 \(60\) 至 \(80\) 分之间；乙部门评分既有 \(30\) 多分，也有 \(100\) 分，分布范围更大. 因此市民对甲部门的总体评价较高且较集中，对乙部门的总体评价较低且较分散.
\end{enumerate}
\end{solution}

\begin{problem}
设 \(F_1\)，\(F_2\) 分别是椭圆 \(C:\frac{x^2}{a^2}+\frac{y^2}{b^2}=1\)（\(a>b>0\)）的左，右焦点，\(M\) 是 \(C\) 上一点且 \(MF_2\) 与 \(x\) 轴垂直，直线 \(MF_1\) 与 \(C\) 的另一个交点为 \(N\).
\begin{enumerate}
\item 若直线 \(MN\) 的斜率为 \(\frac{3}{4}\)，求 \(C\) 的离心率；
\item 若直线 \(MN\) 在 \(y\) 轴上的截距为 \(2\)，且 \(|MN|=5|F_1N|\)，求 \(a\)，\(b\).
\end{enumerate}
\end{problem}

\begin{solution}
\begin{enumerate}
\item 设椭圆焦距的一半为 \(c=\sqrt{a^2-b^2}\)，则 \(F_1=(-c,0)\)，\(F_2=(c,0)\). 因为 \(MF_2\) 垂直于 \(x\) 轴，所以 \(M\) 的横坐标为 \(c\). 将 \(x=c\) 代入椭圆方程，得 \(y_M^2=b^2(1-\frac{c^2}{a^2})=\frac{b^4}{a^2}\). 题设斜率为正，故取 \(M\) 在上半平面，即 \(M=(c,\frac{b^2}{a})\). 直线 \(MN\) 与 \(MF_1\) 是同一直线，故其斜率为 \(\frac{b^2}{2ac}\). 由题设
\(\frac{b^2}{2ac}=\frac34\). 令 \(e=\frac ca\)，并用 \(b^2=a^2-c^2\)，得 \(1-e^2=\frac32e\)，即 \(2e^2+3e-2=0\). 因 \(0<e<1\)，所以 \(e=\frac12\).
\item 直线 \(F_1M\) 在 \(y\) 轴上的截距为 \(2\)，且 \(F_1=(-c,0)\)，故该直线斜率为 \(\frac2c\). 当 \(x=c\) 时 \(y=4\)，所以 \(M=(c,4)\)，从椭圆方程得 \(\frac{c^2}{a^2}+\frac{16}{b^2}=1\)，即 \(b^2=4a\). 在直线上取参数表示 \(X(s)=(-c,0)+s(c,2)\)，则 \(F_1\) 对应 \(s=0\)，\(M\) 对应 \(s=2\). 因 \(F_1\) 在椭圆内部，另一个交点 \(N\) 对应的参数 \(s_N<0\). 由 \(|MN|=5|F_1N|\) 得 \(2-s_N=5(-s_N)\)，所以 \(s_N=-\frac12\)，从而 \(N=(-\frac{3c}{2},-1)\).
又 \(c^2=a^2-b^2=a^2-4a\)，将 \(N\) 代入椭圆方程，得
\(\frac{9c^2}{4a^2}+\frac{1}{b^2}=1\)，即 \(\frac{9(a^2-4a)}{4a^2}+\frac{1}{4a}=1\). 两边乘以 \(4a^2\)，得 \(9a^2-36a+a=4a^2\)，即 \(5a(a-7)=0\). 由 \(a>0\) 得 \(a=7\)，于是 \(b^2=4a=28\)，所以 \(b=2\sqrt7\).
\end{enumerate}
\end{solution}

\begin{problem}
已知函数 \(f(x)=x^3-3x^2+ax+2\)，曲线 \(y=f(x)\) 在点 \((0,2)\) 处的切线与 \(x\) 轴交点的横坐标为 \(-2\).
\begin{enumerate}
\item 求 \(a\)；
\item 证明：当 \(k<1\) 时，曲线 \(y=f(x)\) 与直线 \(y=kx-2\) 只有一个交点.
\end{enumerate}
\end{problem}

\begin{solution}
\begin{enumerate}
\item \(f'(x)=3x^2-6x+a\)，所以曲线在 \((0,2)\) 处的切线方程为 \(y-2=ax\). 它与 \(x\) 轴交点的横坐标为 \(-2\)，故 \(0-2=-2a\)，从而 \(a=1\).
\item 代入 \(a=1\)，交点的横坐标满足
\(x^3-3x^2+x+2=kx-2\)，即 \(g(x)=x^3-3x^2+(1-k)x+4=0\).
当 \(k\leqslant-2\) 时，\(g'(x)=3(x-1)^2-(k+2)\geqslant0\)，且至多在孤立点 \(x=1\) 处为零，故 \(g\) 严格递增. 又 \(g\) 从负无穷趋于正无穷，所以 \(g(x)=0\) 只有一个实根.
当 \(-2<k<1\) 时，令 \(r=\sqrt{\frac{k+2}{3}}\)，则 \(0<r<1\)，且 \(k=3r^2-2\). 此时 \(g'(x)=0\) 的两个根为 \(1-r\)、\(1+r\)，其中 \(1+r\) 是极小值点. 将 \(x=1+s\) 代入并利用 \(k=3r^2-2\)，得 \(g(1+s)=5-3r^2-3r^2s+s^3\). 所以
\(g(1+r)=5-3r^2-2r^3=(1-r)(2r^2+5r+5)>0\).
极小值仍为正，且 \(g(x)\to-\infty\)（\(x\to-\infty\)），\(g(x)\to+\infty\)（\(x\to+\infty\)），故方程只有极小值点左侧的一个实根. 两种情形均表明交点只有一个.
\end{enumerate}
\end{solution}

\begin{problem}
如图，\(P\) 是 \(\odot O\) 外一点，\(PA\) 是切线，\(A\) 为切点，割线 \(PBC\) 与 \(\odot O\) 相交于点 \(B\)，\(C\)，\(PC=2PA\)，\(D\) 为 \(PC\) 的中点，\(AD\) 的延长线交 \(\odot O\) 于点 \(E\). 证明：
\begin{enumerate}
\item \(BE=EC\)；
\item \(AD\cdot DE=2PB^2\).
\end{enumerate}

\bitmapfigure[width=0.40\linewidth]{img/2014/new_curriculum_paper_2_liberal/q22_fig1.png}
\end{problem}

\begin{solution}
\begin{enumerate}
\item 设 \(PA=p\). 由点 \(P\) 的切割线定理，\(PA^2=PB\cdot PC\). 因 \(PC=2p\)，得 \(PB=\frac p2\). 又 \(PD=\frac{PC}{2}=p\)，所以 \(BD=PD-PB=\frac p2=PB\)，\(DC=\frac{PC}{2}=p\)，从而 \(\frac{BD}{DC}=\frac12\).
设 \(\theta=\angle APB=\angle APC\). 在三角形 \(PAB\)、\(PAC\) 中使用余弦定理，得
\(AB^2=p^2+\frac{p^2}{4}-p^2\cos\theta\)，
\(AC^2=p^2+4p^2-4p^2\cos\theta=4AB^2\).
所以 \(AC=2AB\). 因 \(\frac{AB}{AC}=\frac{BD}{DC}\)，由角平分线定理，\(AD\) 平分 \(\angle BAC\). 点 \(E\) 在直线 \(AD\) 上，故 \(\angle BAE=\angle EAC\). 这两个圆周角分别对着弦 \(BE\)、\(EC\)，所以 \(BE=EC\).
\item 因 \(D\) 是弦 \(BC\) 与弦 \(AE\) 的交点，由相交弦定理，\(DB\cdot DC=DA\cdot DE\). 又 \(DB=\frac p2\)，\(DC=p\)，所以
\(AD\cdot DE=DB\cdot DC=\frac{p^2}{2}=2\left(\frac p2\right)^2=2PB^2\).
\end{enumerate}
\end{solution}

\begin{problem}
在直角坐标系 \(xOy\) 中，以坐标原点为极点，\(x\) 轴正半轴为极轴建立极坐标系，半圆 \(C\) 的极坐标方程为 \(\rho=2\cos\theta\)，\(\theta\in\left[0,\frac{\pi}{2}\right]\).
\begin{enumerate}
\item 求 \(C\) 的参数方程；
\item 设点 \(D\) 在 \(C\) 上，\(C\) 在 \(D\) 处的切线与直线 \(l:y=\sqrt{3}x+2\) 垂直，根据（1）中你得到的参数方程，确定点 \(D\) 的坐标.
\end{enumerate}
\end{problem}

\begin{solution}
\begin{enumerate}
\item 由 \(x=\rho\cos\theta\)，\(y=\rho\sin\theta\) 及 \(\rho=2\cos\theta\)，得 \(x=2\cos^2\theta=1+\cos2\theta\)，\(y=2\sin\theta\cos\theta=\sin2\theta\). 令 \(t=2\theta\)，因 \(0\leqslant\theta\leqslant\frac{\pi}{2}\)，所以 \(0\leqslant t\leqslant\pi\). 因此 \(C\) 的参数方程为 \(x=1+\cos t\)，\(y=\sin t\)，\(0\leqslant t\leqslant\pi\).
\item 对参数方程求导，得 \(\frac{dy}{dx}=\frac{\cos t}{-\sin t}=-\cot t\). 直线 \(l\) 的斜率为 \(\sqrt3\)，所以与它垂直的切线斜率为 \(-\frac1{\sqrt3}\). 因而 \(-\cot t=-\frac1{\sqrt3}\)，即 \(\tan t=\sqrt3\). 在 \(0\leqslant t\leqslant\pi\) 内，\(t=\frac{\pi}{3}\). 代入参数方程，得 \(D=(1+\cos\frac{\pi}{3},\sin\frac{\pi}{3})=(\frac32,\frac{\sqrt3}{2})\).
\end{enumerate}
\end{solution}

\begin{problem}
设函数 \(f(x)=\left|x+\frac{1}{a}\right|+\left|x-a\right|\)（\(a>0\)）.
\begin{enumerate}
\item 证明：\(f(x)\geqslant 2\)；
\item 若 \(f(3)<5\)，求 \(a\) 的取值范围.
\end{enumerate}
\end{problem}

\begin{solution}
\begin{enumerate}
\item 由三角不等式和基本不等式，得
\(f(x)=\left|x+\frac1a\right|+|x-a|\geqslant\left|\left(x+\frac1a\right)-(x-a)\right|=a+\frac1a\geqslant2\).
\item 因 \(a>0\)，且 \(3+\frac1a>0\)，分 \(0<a\leqslant3\) 与 \(a>3\) 讨论.
当 \(0<a\leqslant3\) 时，\(f(3)=3+\frac1a+3-a=6+\frac1a-a<5\)，所以 \(a^2-a-1>0\)，结合 \(a>0\) 得 \(a>\frac{1+\sqrt5}{2}\).
当 \(a>3\) 时，\(f(3)=3+\frac1a+a-3=a+\frac1a<5\)，所以 \(a^2-5a+1<0\)，即 \(\frac{5-\sqrt{21}}2<a<\frac{5+\sqrt{21}}2\). 与 \(a>3\) 合并，得到 \(3<a<\frac{5+\sqrt{21}}2\).
综合两种情形，\(a\) 的取值范围为 \(\displaystyle\frac{1+\sqrt5}{2}<a<\frac{5+\sqrt{21}}2\).
\end{enumerate}
\end{solution}
